\begin{textsample}{NEG Dim 2 – vem\_ed\_33 – Score 2.00 – vem\_ed\_33/t173.txt}  \label{ex:f2_neg_045}
Aos Leitores Osvaldo Succi Jr. - Coordenador de Projetos Nesta edição, destacamos a participação da equipe de Apoio à Internacionalização do Ensino Superior (AIES) da Divisão de Extensão e Pesquisa no Ensino Superior (Depes) da Coordenadoria de Ensino Superior de Graduação (CGESG) do Centro Paula Souza (CPS) em eventos online realizados entre novembro de 2025 e fevereiro de 2026: 3ª Jornada COIL UNAM (México); VIII Seminário de Boas Práticas de Ensino e Aprendizagem (SBPEA) da Escola de Engenharia de Lorena da Universidade de São Paulo (EEL-USP) e 18ª Semana de Práticas e Atualizações Pedagógicas (SPAP).
Em 2025, ocorreram as primeiras colaborações entre as Fatecs São José dos Campos e Zona Leste com a Universidade Técnica de Košice, na Eslováquia, envolvendo estudantes brasileiros e eslovacos em projetos de empreendedorismo e desenvolvimento intercultural.
O ano também foi marcado por novos Projetos Colaborativos Internacionais (PCIs) com Portugal, realizados com o Instituto Politécnico de Viseu e o ISP Gaya, envolvendo temáticas como branding, negociação intercultural e modelos preditivos para segurança de redes.
O Artigo de Opinião revela como a gastronomia uniu estudantes do Brasil e da Colômbia em uma rica experiência de Internacionalização em Casa.
Deseja contribuir com Artigo de Opinião?
Escreva para cgesg.pci@cps.sp.gov.br.
Boa leitura!

% matched lemmas: 
\end{textsample}
